\documentclass[tikz,border=10pt]{standalone}
\usepackage[no-math]{fontspec}
\usepackage[brazil]{babel}
\usepackage{amsmath,amssymb}
\IfFileExists{C:/Windows/Fonts/Lato_Regular.ttf}{%
  \setmainfont{Lato}[
    Path = C:/Windows/Fonts/,
    UprightFont = Lato_Regular.ttf,
    BoldFont = Lato_Bold.ttf,
    ItalicFont = Lato_Italic.ttf,
    BoldItalicFont = Lato_BoldItalic.ttf
  ]
  \setsansfont{Lato}[
    Path = C:/Windows/Fonts/,
    UprightFont = Lato_Regular.ttf,
    BoldFont = Lato_Bold.ttf,
    ItalicFont = Lato_Italic.ttf,
    BoldItalicFont = Lato_BoldItalic.ttf
  ]
}{%
  \RequirePackage[default]{lato}
}
\renewcommand{\familydefault}{\sfdefault}

\usepackage{xcolor}
\usepackage{tikz}
\usepackage{pgfplots}
\pgfplotsset{compat=1.18}
\usetikzlibrary{arrows.meta,positioning,shapes.geometric,fit,backgrounds,calc,decorations.pathreplacing,patterns}

% Paleta institucional UEFS / Bioengenharia de Solos
\definecolor{bsNavy}{HTML}{1F4268}
\definecolor{bsNavyDark}{HTML}{152E4A}
\definecolor{bsGold}{HTML}{B38A3D}
\definecolor{bsGoldDark}{HTML}{8C6927}
\definecolor{bsTeal}{HTML}{0E7C7B}
\definecolor{bsForest}{HTML}{2D6A4F}
\definecolor{bsAlert}{HTML}{9E2A2B}
\definecolor{bsEarth}{HTML}{8C5A3C}
\definecolor{bsWater}{HTML}{2E86AB}
\definecolor{bsInk}{HTML}{26313D}
\definecolor{bsMuted}{HTML}{5B6573}

\definecolor{bgNavy}{HTML}{EDF2F7}
\definecolor{bgGold}{HTML}{F8F3E9}
\definecolor{bgTeal}{HTML}{EBF6F5}
\definecolor{bgForest}{HTML}{EEF5F0}
\definecolor{bgAlert}{HTML}{FDF2F2}

\tikzset{
  >=Stealth,
  bsBox/.style={
    draw=#1, fill=white, rounded corners=5pt, line width=1.15pt,
    inner sep=7pt, align=center, font=\footnotesize\sffamily, text=bsInk
  },
  bsTint/.style={
    draw=#1, fill=#1!8!white, rounded corners=5pt, line width=1.15pt,
    inner sep=7pt, align=center, font=\footnotesize\sffamily, text=bsInk
  },
  bsBadge/.style={
    fill=#1, text=white, rounded corners=9pt, inner xsep=8pt, inner ysep=3pt,
    font=\scriptsize\bfseries\sffamily
  },
  bsTitle/.style={font=\normalsize\bfseries\sffamily, text=bsNavy, align=center},
  bsSub/.style={font=\footnotesize\sffamily, text=bsMuted, align=center},
  bsArrow/.style={->, >=Stealth, line width=1.35pt},
  bsDArrow/.style={<->, >=Stealth, line width=1.35pt}
}


\begin{document}
\begin{tikzpicture}

% Título
\node[bsTitle] (title) at (7.3, 5.55)
  {Feixes Vivos e Drenos Verdes: Estrutura Funcional e Disposição em Espinha-de-Peixe};
\node[bsSub] (sub) at (7.3, 5.05)
  {Drenagem subsuperficial combinada ao enraizamento vivo para rebaixamento da poropressão em taludes};

% ================= PAINEL ESQUERDO: SEÇÃO TRANSVERSAL =================
\begin{scope}[shift={(0,0)}]
  \draw[draw=bsForest, fill=bgForest!45, rounded corners=6pt, line width=1.2pt]
    (-0.4, -0.75) rectangle (6.9, 4.40);
  \node[bsBadge=bsForest] at (3.25, 4.05) {SEÇÃO TRANSVERSAL — FEIXE VIVO EM VALA};

  % Solo circundante
  \fill[bsEarth!16] (0.15, 0.30) -- (6.35, 0.30) -- (6.35, 2.75) -- (4.55, 2.75)
    -- (4.10, 1.05) -- (2.40, 1.05) -- (1.95, 2.75) -- (0.15, 2.75) -- cycle;

  % Paredes da vala (0,4 - 0,6 m)
  \draw[bsEarth, line width=1.2pt]
    (0.15, 2.75) -- (1.95, 2.75) -- (2.40, 1.05) -- (4.10, 1.05) -- (4.55, 2.75) -- (6.35, 2.75);

  % Reaterro permeável acima do feixe
  \fill[bsGold!25] (2.15, 2.05) -- (4.35, 2.05) -- (4.52, 2.72) -- (1.98, 2.72) -- cycle;
  \draw[bsGoldDark, dashed, line width=0.8pt] (1.98, 2.72) -- (4.52, 2.72);

  % Feixe cilíndrico de estacas vivas entrelaçadas
  \fill[bsForest!25, draw=bsForest, line width=1.2pt] (3.25, 1.62) ellipse (0.72cm and 0.48cm);
  \foreach \x/\y in {2.85/1.62, 3.12/1.78, 3.38/1.78, 3.65/1.62, 3.00/1.45, 3.28/1.55, 3.52/1.46, 3.25/1.88} {
    \fill[bsEarth!75, draw=bsForest!90!black, line width=0.6pt] (\x, \y) circle (0.11cm);
  }
  \draw[bsNavyDark, line width=1.3pt] (2.53, 1.62) arc (180:360:0.72cm and 0.48cm);

  % Brotação aérea
  \draw[bsForest, line width=1.4pt] (3.05, 2.05) -- (2.75, 3.15);
  \draw[bsForest, line width=1.4pt] (3.25, 2.10) -- (3.25, 3.35);
  \draw[bsForest, line width=1.4pt] (3.45, 2.05) -- (3.78, 3.12);
  \foreach \lx/\ly in {2.75/3.15, 2.90/2.85, 3.25/3.35, 3.12/2.95, 3.78/3.12, 3.62/2.80} {
    \fill[bsForest!85] (\lx, \ly) ellipse (0.13cm and 0.06cm);
  }

  % Raízes para ancoragem no solo
  \draw[bsEarth!90!black, line width=0.9pt] (2.85, 1.22) .. controls (2.55, 0.85) .. (2.15, 0.55);
  \draw[bsEarth!90!black, line width=0.9pt] (3.25, 1.14) .. controls (3.20, 0.78) .. (3.05, 0.45);
  \draw[bsEarth!90!black, line width=0.9pt] (3.65, 1.22) .. controls (3.95, 0.85) .. (4.35, 0.55);

  % Setas de drenagem subsuperficial lateral
  \draw[bsArrow, bsWater, line width=1.4pt] (0.45, 1.65) -- (2.05, 1.65);
  \draw[bsArrow, bsWater, line width=1.4pt] (6.05, 1.65) -- (4.45, 1.65);

  % Rótulos claros sem sobreposição
  \node[font=\scriptsize\bfseries\sffamily, text=bsWater, align=center] at (1.10, 2.18)
    {Drenagem\\subsuperficial};
  \node[font=\scriptsize\bfseries\sffamily, text=bsWater, align=center] at (5.40, 2.18)
    {Fluxo\\lateral};
  \node[font=\scriptsize\bfseries\sffamily, text=bsGoldDark, align=center] at (5.35, 3.15)
    {Reaterro\\permeável};
  \draw[->, bsGoldDark, line width=0.8pt] (4.70, 3.00) -- (4.05, 2.45);

  \node[font=\scriptsize\bfseries\sffamily, text=bsForest, align=center] at (1.25, 3.30)
    {Estacas vivas\\entrelaçadas};
  \draw[->, bsForest, line width=0.8pt] (2.05, 3.22) -- (2.82, 2.95);

  \node[font=\scriptsize\bfseries\sffamily, text=bsEarth!90!black, align=center] at (1.25, 0.68)
    {Vala $0{,}4\text{--}0{,}6\text{ m}$};
  \node[font=\scriptsize\bfseries\sffamily, text=bsEarth!90!black, align=center] at (5.20, 0.68)
    {Ancoragem\\radicular};

  \node[fill=white, draw=bsForest, rounded corners=4pt, text width=6.4cm, inner sep=4pt,
        font=\scriptsize\sffamily, text=bsInk, align=center] at (3.25, -0.22)
    {\textbf{Função:} Alívio imediato da poropressão + reforço mecânico};
\end{scope}

% ================= PAINEL DIREITO: PLANTA ESPINHA-DE-PEIXE =================
\begin{scope}[shift={(7.7,0)}]
  \draw[draw=bsTeal, fill=bgTeal!45, rounded corners=6pt, line width=1.2pt]
    (-0.4, -0.75) rectangle (6.9, 4.40);
  \node[bsBadge=bsTeal] at (3.25, 4.05) {PLANTA — DISPOSIÇÃO EM ESPINHA-DE-PEIXE};

  % Face do talude saturado
  \fill[bsEarth!10, draw=bsMuted!50, rounded corners=4pt] (0.35, 0.62) rectangle (6.15, 3.58);

  % Indicação de declividade
  \draw[bsArrow, bsMuted, line width=1.1pt] (0.85, 3.28) -- (0.85, 1.40);
  \node[font=\scriptsize\bfseries\sffamily, text=bsMuted, align=center, rotate=90]
    at (0.58, 2.34) {Sentido do declive};

  % Dreno principal (coletor central)
  \draw[bsForest, line width=4.2pt, rounded corners=2pt] (3.35, 3.38) -- (3.35, 0.72);
  \draw[bsArrow, bsWater, line width=1.3pt] (3.35, 1.38) -- (3.35, 0.74);

  % Drenos secundários diagonais simétricos (3 pares)
  \foreach \y in {3.10, 2.40, 1.70} {
    \draw[bsTeal, line width=2.8pt] (1.65, \y+0.35) -- (3.35, \y-0.30);
    \draw[bsTeal, line width=2.8pt] (5.05, \y+0.35) -- (3.35, \y-0.30);
    \draw[bsArrow, bsWater, line width=0.95pt] (2.05, \y+0.20) -- (2.75, \y-0.07);
    \draw[bsArrow, bsWater, line width=0.95pt] (4.65, \y+0.20) -- (3.95, \y-0.07);
  }

  % Exutório estabilizado na base
  \node[fill=bsNavy, text=white, rounded corners=4pt, inner xsep=7pt, inner ysep=3pt,
        font=\scriptsize\bfseries\sffamily] at (3.35, 0.45)
    {Exutório Estabilizado / Dissipador};

  % Rótulos internos claros em espaço livre
  \node[fill=white, draw=bsTeal, rounded corners=3pt, inner xsep=5pt, inner ysep=2.5pt,
        font=\scriptsize\bfseries\sffamily, text=bsTeal, align=center] at (5.38, 1.25)
    {Drenos\\Secundários};
  \draw[->, bsTeal, line width=0.8pt] (4.95, 1.55) -- (4.50, 1.82);

  \node[fill=white, draw=bsForest, rounded corners=3pt, inner xsep=5pt, inner ysep=2.5pt,
        font=\scriptsize\bfseries\sffamily, text=bsForest, align=center] at (1.90, 1.10)
    {Dreno\\Principal};
  \draw[->, bsForest, line width=0.8pt] (2.38, 1.20) -- (3.22, 1.48);

  \node[fill=white, draw=bsTeal, rounded corners=4pt, text width=6.4cm, inner sep=4pt,
        font=\scriptsize\sffamily, text=bsInk, align=center] at (3.25, -0.22)
    {\textbf{Rede:} Captação simétrica até exutório protegido};
\end{scope}

\end{tikzpicture}
\end{document}
